\section{Back-office : un effet \enquote{boîte noire}}
Malgré la robustesse du logiciel XTF, l'un des problèmes majeurs de celui-ci vient du manque d'autonomie des équipes métiers face à cette infrastructure fermée et non documentée. Cela rend le travail interne davantage plus complexe notamment lorsque l'un d'eux repèrent une erreur syntaxique par exemple, ils peuvent la corriger mais ensuite ils doivent relancer toute la procédure de publication. A l'origine, où une petite correction peut prendre entre une et deux minutes à être corrigé, pour le moment ce type d'erreur syntaxique ou une erreur plus importante mettent minimum trois jours à être corrigé et à apparaître dans la bibliothèque numérique pour être consulté par l'usager. En effet, une fois l'erreur rectifiée, la personne en charge de la publication doit à nouveau rouvrir le PDF qu'elle avait scellé un mot de passe afin que celui-ci ne soit pas modifiable pour un lecteur. L'expert métier déverouille le PDF, modifie la date de création du PDF, vérouille à nouveau le PDF et continue la procédure de mise en ligne. De plus, à chaque étape, ils doivent attendre les traitements de nuits valident les étapes précédentes avant de pouvoir poursuivre les modifications, ajoutant une lourdeur à la procédure de travail. 

Le second problème vient d'XTF lui-même qui n'est plus maintenu. Habituellement, les inconvénients quotidiens ou difficultés rencontrées par les experts métiers étaient tous palier par le responsable du département informatique technique (DIT). Le risque n'est pour autant pas moindre car tous les services de la bibliothèque numérique actuelle qu'il soit interne ou externe repose uniquement sur les compétences d'une seule et même personne. Aujourd'hui, la bibliothèque numérique actuelle survit uniquement grâce aux compétences informatiques du responsable qui a su trouver des solutions comme avec l'ajout en 2013 d'une visionneuse JPEG apportant un confort pour la lecture  des PDF notamment ceux volumineux. 

La situation de la bibliothèque Cujas est loin d'être un cas isolé celle-ci illustre bien la pensée de l'époque des institutions patrimoniales. Le cas de Cujas illustre le concept du deuxième âge théorisé par Emmanuelle Bermès et Gautier Poupeau. Ce deuxième âge s'étend des années 1990 à 2010. Cette période est marquée par deux événements majeurs : l'irruption du Web au milieu des années 1990 et la numérisation de masse insufflé par Google et Microsoft à partir des années 2000. C'est la période correspondant aux développements des projets de numérisations de la bibliothèque Cujas qui s'est structurée selon les pratiques de travail de l'époque. Par ailleurs, c'est à ce moment-là qu'émerge les premières bibliothèques numériques. 

Entre 1990 et 2010, les institutions patrimoniales ont une vision différente de la conception d'un document numérique. A l'épopque, les bibliothèques voient le document numérique comme une conversion du document physique. Jean-Michel Salaün théorise d'ailleurs cette idée qu'il définit \enquote{comme un ensemble de données organisées selon une structure fiable associées à des règles de mise en forme permettant une lisibilité partagée entre son concepteur et ses lecteurs}. Avec l'émerge du web, on voit se développer dans le paysage des technologies numériques des standards particulièrement utile pour structurée l'ouvrage. C'est l'arrivée massif du format XML qui permet le triomphe \enquote{de la pensée hiérarchique et documentaire}, modélisant les documents sous forme d'arborescences. Le format XML est alors pensé pour permettre la diffusion des données avec les autres bibliothèques sauf qu'avec l'arrivée des ordinateurs personnels freinent paradoxamelement l'innovation. Christine Ducourtieux l'exprime d'ailleurs ainsi : \enquote{l'arrivée des systèmes d'exploitation Windows et MacIntosh et de la bureautique ont paradomalement pour effet de freiner l'élan d'innovation technique dans les méthodes, remplacé progressivement par des patriques discrètes : une appropriation de l'outil numérique plus individuelle, qui n'est plus subordonnée à l'histoire quantitative}. Le développement de la technologie et l'arrivée de l'Internet en 1989, entraîne à nouveau une forme de cloisonnement avec un retour à l'individualisation des pratiques des chercheurs. 

En effet, l'arrivée des ordinateurs personnels bouleversent les pratiques créant une \enquote{séparation nette} entre l'édition scientifique et la numérisation de masse. Emmanuelle Bermès justifie cette hypothèse en expliquant que \enquote{l'édition scientifique numérique, qui reste le domaine de la recherche et peut s'adapter aux spécificités des corpus, et la numérisation, qui vise des masses documentaires plus importantes et traitées de manière homogène}. Au lieu de travailler ensemble, le développement informatique amène à revenir à des pratiques individuelles. Les bibliothèques numériques sont alors uniquement pensées comme des catalogues de PDF n'ayant nullement besoin de valorisation ne s'appuyant donc pas sur le travail des chercheurs. 

On observe aussi une parcellisation industrielle des tâches en bibliothèque. Il s'agit d'une époque où l'on segmentise énormément le travail de chaque personne sans lui donner accès au reste de la chaîne de production considérant que cela n'est pas dans ces attributions et ne lui est donc pas nécessaire pour réaliser son travail. La bibliothèque Cujas est héritière de cette manière de travail mais c'est loin d'être un cas isolé d'autres institutions patrimoniales comme la BnF ont également rencontré ce problème. Michel Richard écrit sur ce sujet :\enquote{l'équipe responsable de la réalisation de ce programme comprend quinze personnes à plein temps rattachées sur le site d'Ivry. elle a en charge l'ensemble des tâches nécessaires à la constitution de la collection (travail bibliographique, contacts avece les universités et l'édition, relations avec la recherche, gestion du fonds électronique, commandes aux libraires, négocations avec les ayants droit, etc.) et au suiv de production (choix des prestataires, tests, contrôle qualité, gestion des stocks, normalisation, etc.)}. Cela alourdit énormément le travail des bibliothécaires car c'est uniquement le \enquote{résultat final} qui est mutualisé et non le reste.  

La bibliothèque Cujas fait d'ailleurs face à cette contrainte en disant clairement que le système actuel survit uniquement grâce aux compétences informatiques du responsable du DIT. En dehors de lui, personne n'est à même à Cujas de pouvoir régler le problème entraînant ainsi une centralisation des compétences en une seule personne mettant en péril la pérennité de la bibliothèque numérique dans le futur. Les données deviennent alors vulnérable dans un monde où avec l'arrivée du Web sémantique, on réfléchit à ouvrir d'avantage les données afin de permettre une libre circulation de celle-ci.